\documentclass[12pt,letterpaper]{article}
\usepackage[utf8]{inputenc}
\usepackage[T1]{fontenc}
\usepackage[spanish,es-nodecimaldot]{babel}
\usepackage{newtxtext}
\usepackage{newtxmath}
\usepackage[letterpaper,top=2.5cm,bottom=2.5cm,left=3cm,right=3cm,footskip=1.25cm]{geometry}
\usepackage{setspace}
\usepackage{ragged2e}
\usepackage{fancyhdr}
\usepackage{titlesec}
\usepackage{enumitem}
\usepackage{needspace}
\usepackage{microtype}
\usepackage[hidelinks]{hyperref}
\hypersetup{pdftitle={Propuesta de capitulado y justificación},pdfauthor={Carlos Emiliano Segura Loera, Valeria Hernandez Santos, Nehiel Ramirez Hernandez},pdfsubject={Investigación en Matemáticas Aplicadas y Computación}}
\onehalfspacing
\setlength{\parindent}{1.25cm}
\setlength{\parskip}{4pt}
\pagestyle{fancy}
\fancyhf{}
\fancyfoot[L]{\small FES Acatlán · UNAM}
\fancyfoot[R]{\small\thepage}
\renewcommand{\headrulewidth}{0pt}
\renewcommand{\footrulewidth}{0.4pt}
\titleformat{\section}{\normalfont\bfseries\large}{\thesection.}{0.6em}{}
\titleformat{\subsection}{\normalfont\bfseries\normalsize}{\thesubsection.}{0.6em}{}
\titlespacing*{\section}{0pt}{16pt}{9pt}
\titlespacing*{\subsection}{0pt}{12pt}{6pt}
\raggedbottom
\newcommand{\numerocap}{1}
\newlist{apartados}{enumerate}{2}
\setlist[apartados,1]{label=\numerocap.\arabic*.,leftmargin=1.05cm,labelsep=0.18cm,itemsep=2pt,topsep=4pt,parsep=0pt}
\setlist[apartados,2]{label=\numerocap.\arabic{apartadosi}.\arabic*.,leftmargin=1.35cm,labelsep=0.18cm,itemsep=1pt,topsep=2pt,parsep=0pt}
\newcommand{\capitulo}[2]{\renewcommand{\numerocap}{#1}\Needspace{5\baselineskip}\subsection*{Capítulo #1. #2}}
% Fecha de elaboración de esta actividad; modificar si se indica otra entrega.
\newcommand{\fechaactividad}{18 de octubre de 2026}
\begin{document}
\begin{titlepage}
\thispagestyle{empty}
\begin{spacing}{1}
\begin{center}
\vspace*{0.5cm}
{\large\bfseries UNIVERSIDAD NACIONAL AUTÓNOMA DE MÉXICO\par}
\vspace{0.25cm}
{\large FACULTAD DE ESTUDIOS SUPERIORES ACATLÁN\par}
\vspace{0.15cm}
{\normalsize Licenciatura en Matemáticas Aplicadas y Computación\par}
\vspace{1.15cm}
{\large\bfseries INVESTIGACIÓN EN MATEMÁTICAS APLICADAS Y\\COMPUTACIÓN\par}
\vspace{1.1cm}
{\LARGE\bfseries Contenido sintético en Internet y degradación\\[0.15cm]de modelos de lenguaje\par}
\vspace{0.35cm}
{\large Una aproximación al fenómeno de \textit{Model Collapse}\par}
\vspace{0.7cm}
{\Large\bfseries Propuesta de capitulado\par}
\vspace{0.2cm}
{\normalsize Estructura provisional y justificación\par}
\end{center}
\vspace{1.4cm}
\begin{flushright}
\begin{minipage}{0.64\textwidth}
\textbf{Alumnos:}\\[0.15cm]
Carlos Emiliano Segura Loera\\
Valeria Hernandez Santos\\
Nehiel Ramirez Hernandez
\vspace{0.5cm}

\textbf{Profesor:}\\[0.1cm]
Sergio Quiroz Almaraz
\vspace{0.5cm}

\textbf{Fecha de elaboración:}\\[0.1cm]
\fechaactividad
\vspace{0.5cm}

Naucalpan de Juárez, Estado de México
\end{minipage}
\end{flushright}
\end{spacing}
\end{titlepage}
\pagenumbering{arabic}
\justifying
\section*{Presentación de la propuesta}
El presente documento propone el capitulado de la investigación
\textit{Contenido sintético en Internet y degradación de modelos de lenguaje}.
La organización parte de la siguiente pregunta principal:
\begin{quote}
¿Cómo cambia la diversidad del contenido generado por modelos de lenguaje
cuando se incorporan repetidamente datos sintéticos al proceso de entrenamiento?
\end{quote}
Se plantea una estructura de siete capítulos centrales y un capítulo
complementario opcional que permita avanzar del problema y sus fundamentos
hacia el diseño experimental, la presentación de resultados y su
interpretación. La proporción de datos sintéticos y el número de
generaciones organizan la comparación; la diversidad y la repetición
constituyen las principales dimensiones de evaluación.

El capitulado aún puede cambiar. Los apartados indican qué se propone
desarrollar, pues aún no se han definido por completo el modelo, el conjunto
de datos ni las métricas que se utilizarán.

\section*{Capitulado propuesto}
\begin{spacing}{1.08}
\capitulo{1}{Introducción}
\begin{apartados}
\item Contexto del contenido sintético y los modelos de lenguaje
\item Motivación del estudio
\begin{apartados}
\item La teoría del Internet muerto como contexto
\item Reutilización de contenido sintético en el entrenamiento
\end{apartados}
\item Planteamiento del problema
\item Pregunta principal de investigación
\item Preguntas secundarias
\item Objetivos de la investigación
\begin{apartados}
\item Objetivo general
\item Objetivos específicos
\end{apartados}
\item Hipótesis de investigación
\item Justificación y aportación esperada
\item Alcances y delimitaciones
\end{apartados}

\capitulo{2}{Marco teórico}
\begin{apartados}
\item Modelos generativos de lenguaje
\begin{apartados}
\item Generación de texto como proceso probabilístico
\item Preentrenamiento y ajuste de modelos
\item Estrategias de muestreo y decodificación
\end{apartados}
\item Datos utilizados en el entrenamiento
\begin{apartados}
\item Datos de origen humano
\item Datos sintéticos y semisintéticos
\item Mezcla, reemplazo y acumulación de datos
\end{apartados}
\item Entrenamiento recursivo
\begin{apartados}
\item Ciclos de entrenamiento y generación
\item Transmisión de información entre generaciones
\end{apartados}
\item El fenómeno de \textit{model collapse}
\begin{apartados}
\item Definición y manifestaciones
\item Cambios en la distribución del contenido
\item Pérdida de patrones poco frecuentes
\item Diferencias entre colapso, repetición y disminución del rendimiento
\end{apartados}
\item Diversidad del contenido generado
\begin{apartados}
\item Diversidad léxica
\item Diversidad semántica
\item Diversidad sintáctica
\end{apartados}
\item Repetición y concentración de patrones
\begin{apartados}
\item Repetición dentro de un texto
\item Similitud entre textos generados
\end{apartados}
\item Fundamentos de las métricas de evaluación
\begin{apartados}
\item Medidas de diversidad léxica
\item Medidas de diversidad semántica
\item Medidas de repetición
\item Medidas de fidelidad distribucional y rendimiento
\end{apartados}
\end{apartados}

\capitulo{3}{Estado del arte}
\begin{apartados}
\item Criterios de búsqueda y selección de literatura
\item Evidencia inicial del colapso por entrenamiento recursivo
\item Estudios sobre pérdida de diversidad lingüística
\item Efectos de la proporción de datos sintéticos
\item Conservación de datos originales y acumulación de generaciones
\item Influencia de las estrategias de generación
\item Estrategias de mitigación
\begin{apartados}
\item Selección y filtrado de datos
\item Generación de datos semisintéticos
\item Modificación del objetivo de entrenamiento
\end{apartados}
\item Comparación crítica de los trabajos revisados
\begin{apartados}
\item Modelos, datasets y tareas
\item Protocolos y presupuestos de entrenamiento
\item Métricas y limitaciones
\end{apartados}
\item Delimitación de la aportación del presente estudio
\end{apartados}

\capitulo{4}{Metodología}
\begin{apartados}
\item Enfoque y diseño experimental
\item Definición operacional de las variables
\begin{apartados}
\item Proporción de datos sintéticos
\item Número de generaciones
\item Diversidad léxica y semántica
\item Repetición del contenido
\item Variables de control
\end{apartados}
\item Recursos e instrumentos computacionales
\begin{apartados}
\item Hardware disponible
\item Software, bibliotecas y versiones
\item Herramientas de entrenamiento y evaluación
\end{apartados}
\item Selección del modelo de lenguaje
\begin{apartados}
\item Criterios de selección
\item Características y configuración del modelo
\end{apartados}
\item Selección y preparación del conjunto de datos
\begin{apartados}
\item Procedencia, idioma y características
\item Limpieza y eliminación de duplicados
\item División en entrenamiento, validación y prueba
\item Identificación de patrones frecuentes y poco frecuentes
\end{apartados}
\item Condiciones experimentales
\begin{apartados}
\item Condición de referencia sin datos sintéticos
\item Proporciones sintéticas seleccionadas
\item Número de generaciones y repeticiones independientes
\item Política de conservación o reemplazo de datos originales
\item Presupuesto de datos y entrenamiento por condición
\end{apartados}
\item Procedimiento de entrenamiento recursivo
\begin{apartados}
\item Preparación del modelo inicial
\item Generación del contenido sintético
\item Construcción del conjunto de entrenamiento siguiente
\item Inicialización y ajuste del modelo en cada generación
\item Evaluación y almacenamiento de resultados
\end{apartados}
\item Configuración de generación del texto
\begin{apartados}
\item Selección de instrucciones o contextos de entrada
\item Longitud y cantidad de textos generados
\item Parámetros de muestreo y decodificación
\item Gestión de semillas aleatorias
\end{apartados}
\item Métricas e instrumentos de evaluación
\begin{apartados}
\item Cálculo de diversidad léxica
\item Cálculo de diversidad semántica
\item Cálculo de repetición dentro y entre textos
\item Evaluación de patrones poco frecuentes
\item Evaluación complementaria de rendimiento o fidelidad
\end{apartados}
\item Procedimiento de análisis estadístico
\begin{apartados}
\item Unidad de análisis y unidad de réplica experimental
\item Estadísticos descriptivos e incertidumbre
\item Comparación entre proporciones y generaciones
\item Análisis de la interacción entre ambos factores
\item Criterios de contrastación de la hipótesis
\end{apartados}
\item Controles de reproducibilidad y calidad
\begin{apartados}
\item Registro de configuraciones y versiones
\item Prevención de filtración de datos de evaluación
\item Verificación de duplicados y errores del procedimiento
\end{apartados}
\item Consideraciones éticas y licencias de uso
\end{apartados}

\capitulo{5}{Resultados}
\begin{apartados}
\item Ejecución del experimento y datos obtenidos
\item Resultados de la condición de referencia
\item Evolución de la diversidad léxica
\begin{apartados}
\item Diferencias entre generaciones
\item Diferencias entre proporciones sintéticas
\end{apartados}
\item Evolución de la diversidad semántica
\begin{apartados}
\item Diferencias entre generaciones
\item Diferencias entre proporciones sintéticas
\end{apartados}
\item Evolución de la repetición
\begin{apartados}
\item Repetición dentro de los textos
\item Similitud entre textos
\end{apartados}
\item Cambios en la representación de patrones poco frecuentes
\item Resultados de la evaluación complementaria
\item Interacción entre proporción sintética y generaciones
\item Variabilidad entre repeticiones experimentales
\item Resultados de los contrastes estadísticos
\end{apartados}

\capitulo{6}{Discusión}
\begin{apartados}
\item Interpretación de los cambios en diversidad
\item Interpretación de los cambios en repetición
\item Efecto conjunto de la proporción sintética y las generaciones
\item Relación entre diversidad, fidelidad y rendimiento
\item Evidencia compatible con \textit{model collapse}
\begin{apartados}
\item Indicadores observados
\item Explicaciones alternativas y alcance de la evidencia
\end{apartados}
\item Comparación con los estudios previos
\item Evaluación de la hipótesis de investigación
\item Limitaciones y amenazas a la validez
\begin{apartados}
\item Modelo, dataset e idioma
\item Presupuesto computacional y número de repeticiones
\item Limitaciones de las métricas
\item Alcance de la generalización de los resultados
\end{apartados}
\end{apartados}

\capitulo{7}{Conclusiones}
\begin{apartados}
\item Respuesta a la pregunta principal
\item Respuestas a las preguntas secundarias
\item Cumplimiento de los objetivos
\item Conclusión sobre la hipótesis, dentro de las condiciones evaluadas
\item Aportaciones del estudio
\item Limitaciones y líneas de continuación (si no se incluye el capítulo 8)
\end{apartados}

\capitulo{8}{Trabajo futuro (opcional)}
\noindent\textit{La inclusión de este capítulo estará sujeta al tiempo
disponible antes de la entrega. En esta propuesta se considera un apartado
complementario.}
\begin{apartados}
\item Evaluación con otros modelos, tareas e idiomas
\item Ampliación de generaciones y repeticiones
\item Comparación de políticas de conservación de datos
\item Evaluación de estrategias de mitigación
\end{apartados}

\Needspace{6\baselineskip}
\subsection*{Referencias bibliográficas}
Fuentes citadas a lo largo de la investigación, con un formato de referencia
uniforme y datos bibliográficos verificados.

\subsection*{Anexos}
\begin{enumerate}[label=\Alph*.,leftmargin=0.8cm,itemsep=2pt,topsep=4pt]
\item Configuraciones completas del experimento.
\item Instrucciones y contextos utilizados para generar texto.
\item Definiciones y fórmulas de las métricas seleccionadas.
\item Tablas completas de resultados.
\item Muestras de contenido por condición y generación.
\item Código y guía de reproducción.
\end{enumerate}
\end{spacing}

\clearpage
\section*{Justificación de la selección y organización}

\subsection*{Criterio general}
Se selecciona una estructura que vincula la pregunta de investigación con las
variables, los procedimientos y las evidencias que permitirán responderla.
El orden propuesto distingue la explicación conceptual, la revisión de
antecedentes, la ejecución del experimento y la interpretación de sus
resultados. Esta separación facilita desarrollar cada tema con una función
definida y evita repetir contenidos entre capítulos.

La granularidad es mayor en la metodología porque el trabajo requiere
precisar cómo se construyen las generaciones de datos y cómo se comparan.
En cambio, la introducción y las conclusiones se mantienen centradas en el
problema y su respuesta. Los niveles de numeración permiten convertir cada
apartado en una tarea de redacción o de preparación experimental.

\subsection*{Capítulo 1. Introducción}
Se incluyen el problema, las preguntas, los objetivos, la hipótesis y la
delimitación para establecer desde el comienzo qué se investigará y por qué.
La teoría del Internet muerto se incorpora únicamente como motivación, pues
la investigación no busca comprobarla ni estimar la proporción de contenido
sintético en Internet.

El tratamiento técnico de los modelos y del colapso se reserva al marco
teórico. De esta manera, la introducción puede presentar el propósito del
estudio sin extenderse en explicaciones que se desarrollarán posteriormente.

\subsection*{Capítulo 2. Marco teórico}
Se seleccionan los conceptos necesarios para comprender el experimento:
modelos de lenguaje, datos sintéticos, entrenamiento recursivo, distribución,
diversidad y repetición. La distinción entre preentrenamiento y ajuste
permitirá describir con precisión el procedimiento que finalmente se utilice.

La separación entre diversidad léxica, semántica y sintáctica ayuda a evitar
que toda variación del contenido se interprete como un mismo fenómeno.
La diversidad sintáctica se incluye como fundamento conceptual; su medición
experimental dependerá de la factibilidad y no constituye todavía un
compromiso metodológico. También se distingue el colapso de la repetición y
del deterioro del rendimiento para no identificarlo a partir de un único
indicador.

\subsection*{Capítulo 3. Estado del arte}
Este capítulo se organiza por problemas y enfoques, de modo que la literatura
permita comparar cómo se ha estudiado el fenómeno. Se consideran la
proporción sintética, la conservación de datos originales, la acumulación,
las estrategias de generación y las mitigaciones porque son decisiones
directamente relacionadas con el diseño del proyecto.

La comparación de modelos, tareas, métricas y presupuestos permitirá
delimitar una aportación realista. El trabajo no pretende presentar el
\textit{model collapse} como un fenómeno nuevo; busca definir una reproducción
o comparación experimental concreta, accesible y documentada.

\subsection*{Capítulo 4. Metodología}
Se propone el mayor nivel de detalle en este capítulo para que el
procedimiento pueda ejecutarse y reproducirse. Se justificará la selección
de los modelos pequeños, se describirán los datos, se definirán las
condiciones y se registrarán los instrumentos computacionales utilizados.

Se incluyen una condición sin datos sintéticos, la política de conservación
de datos humanos y el presupuesto de entrenamiento para distinguir los
efectos de las condiciones comparadas. También se propone precisar la
inicialización de cada ciclo, los parámetros de generación y las semillas
aleatorias, pues forman parte del procedimiento que produce el contenido
evaluado.

El análisis estadístico se incorpora desde el diseño para precisar las
unidades de réplica, la incertidumbre y los criterios de comparación.
Los textos generados por una misma cadena de entrenamiento no se tratarán
automáticamente como experimentos independientes. Las pruebas y métricas
concretas se seleccionarán después de definir el diseño y los datos
disponibles.

\subsection*{Capítulo 5. Resultados}
Los resultados se organizan según las dimensiones medidas y los factores
experimentales. Esta distribución permite presentar cómo evolucionan la
diversidad y la repetición entre generaciones y proporciones, además de
mostrar la variabilidad entre repeticiones.

Se incluyen la referencia sin datos sintéticos y los contrastes estadísticos
para que las comparaciones tengan un punto de partida explícito. En este
capítulo se reportarán mediciones, tablas, gráficas y ejemplos obtenidos con
criterios previamente definidos; las explicaciones de sus causas y su
relación con la literatura se desarrollarán en la discusión.

\subsection*{Capítulo 6. Discusión}
La discusión se mantiene separada de los resultados para distinguir lo
observado de su interpretación. Aquí se examinará qué significan los cambios,
cómo se relacionan con los estudios previos y qué evidencia aportan respecto
de la hipótesis.

Se incorpora un apartado sobre explicaciones alternativas porque una
disminución de diversidad o un aumento de repetición, por sí solos, pueden
ser insuficientes para sostener una conclusión sobre colapso. Asimismo,
las limitaciones del modelo, los datos, las métricas y el presupuesto
permitirán precisar hasta dónde se pueden generalizar los hallazgos.

\subsection*{Capítulo 7. Conclusiones}
Se seleccionan apartados vinculados directamente con la pregunta principal,
las preguntas secundarias, los objetivos y la hipótesis. Esta correspondencia
evita cerrar el documento con afirmaciones generales sobre inteligencia
artificial que no se deriven del experimento.

Las conclusiones deberán expresar lo que los resultados permitan sostener,
incluida la posibilidad de no observar el efecto esperado o de encontrarlo
solo en determinadas condiciones. La aportación se describirá de acuerdo con
el alcance real del trabajo.

Si no se incluye el capítulo de Trabajo futuro, el apartado 7.6 reunirá
una síntesis de las limitaciones y de las posibles líneas de continuación,
sin repetir su análisis detallado en la discusión.

\subsection*{Capítulo 8. Trabajo futuro (opcional)}
Se propone un capítulo breve para las extensiones que excedan el alcance
actual: otros modelos, idiomas, tareas, más generaciones o estrategias de
mitigación. Su función es identificar preguntas que surjan del estudio sin
convertirlas en compromisos del experimento presente.

Se considera este apartado complementario, pero de interés para dar
continuidad al estudio. Su inclusión dependerá del tiempo disponible.
Se dará prioridad
al análisis y a la respuesta de las preguntas de investigación. Si el tiempo
es limitado, las posibles continuaciones podrán mencionarse brevemente en
las conclusiones, a partir de las limitaciones o los hallazgos del estudio,
sin que su realización forme parte de la entrega actual.

\subsection*{Correspondencia con las preguntas de investigación}
La organización propuesta permite dar seguimiento a las preguntas
secundarias mediante apartados concretos:
\begin{itemize}[leftmargin=0.7cm,itemsep=4pt]
\item \textbf{Diversidad léxica y semántica:} se define en el marco teórico,
se operacionaliza en la metodología y se presenta en los apartados 5.3 y 5.4.
\item \textbf{Proporción de datos sintéticos:} se establece como variable
experimental y se compara dentro de los resultados de cada dimensión.
\item \textbf{Repetición:} se evalúa dentro de cada texto y entre textos,
con resultados concentrados en el apartado 5.5.
\item \textbf{Interacción entre proporción y generaciones:} se incorpora
en el diseño del análisis y se presenta específicamente en el apartado 5.8.
\item \textbf{Indicadores asociados con \textit{model collapse}:} se examinan
en el apartado 6.5 considerando conjuntamente los resultados y las
explicaciones alternativas.
\end{itemize}

\subsection*{Carácter provisional y alcance}
Los apartados de patrones poco frecuentes y evaluación complementaria
dependen de contar con una referencia y un procedimiento de medición
adecuados. Si no resultan viables, deberán ajustarse el capitulado y las
afirmaciones que se pretenda sostener. La estructura propuesta organiza el
desarrollo del trabajo, pero no anticipa resultados ni sustituye la definición
final de la metodología.

\end{document}
